\pdfmapfile{+symbols.map}
\pdfmapfile{+cm.map}
\pdfmapfile{+cmextra.map}
\documentclass[11pt]{article}
\usepackage[margin=1in]{geometry}
\usepackage{amsmath,amssymb,amsthm}
\usepackage[colorlinks=true,urlcolor=blue]{hyperref}
\newtheorem{theorem}{Theorem}
\title{A four cycle counterexample to conjecture 00000000579}
\author{Galaxy-0}
\date{10 October 2026}
\begin{document}
\sloppy
\maketitle
\begin{abstract}
The simple cycle on four vertices disproves the conjectured integer divisibility. Its graphic matroid has Kazhdan--Lusztig polynomial $1+2t$, its dual has polynomial $1$, and the absolute value of its chromatic polynomial at $-1$ is $14$. The difference at $t=1$ is $2$, which is not divisible by $14$. A self-contained Lean 4 project derives the values from the defining recurrence, checks an actual graphic rank function, and verifies the existence of recurrence solutions.
\end{abstract}

\section{The assertion being refuted}
The English statement asserts, for every graphic matroid $M(G)$, that
\[
 P_{M(G)}(1)-P_{M(G)^*}(1)
 \quad\hbox{is an integer multiple of}\quad |\chi_G(-1)|.
\]
The Chinese statement asserts the same integer-multiple condition. The further requirement that the multiplier be a spanning-tree deletion count cannot rescue a failure of integer divisibility. We refute the displayed necessary condition without assigning a meaning to that additional count.

Here $\chi_G$ is the chromatic polynomial. If instead it were read as the characteristic polynomial of the graphic matroid, the same example and the same absolute value apply. Both readings are checked below.

\section{The graph and its two rank functions}
Let $V=\{0,1,2,3\}$ and let edge $i$ join $i$ to $i+1$ modulo $4$. Thus $G=C_4$ is a connected simple graph with four edges. Every proper edge subset is a forest. The rank of its cycle matroid $M$ is
\[
 r(S)=4-c(V,S)=\min(3,|S|),\qquad S\subseteq E,
\]
where $c(V,S)$ counts connected components, including isolated vertices. Therefore $M=U_{3,4}$. The usual dual-rank formula gives
\[
 r^*(S)=|S|+r(E\setminus S)-r(E)=\min(1,|S|),
\]
so $M^*=U_{1,4}$. These are ordinary matroids; no special convention for a degenerate graph is needed.

\section{Deriving the Kazhdan Lusztig polynomials}
For a loopless matroid $N$, its Kazhdan--Lusztig polynomial is characterized by the rank-zero value $1$, the strict bound $\deg P_N<r(N)/2$ at positive rank, and the recurrence
\begin{equation}\label{rec}
 t^{r(N)}P_N(t^{-1})=\sum_{F\text{ flat of }N}\chi_{N|F}(t)P_{N/F}(t).
\end{equation}
This is the defining characterization of Elias, Proudfoot and Wakefield, Theorem 2.2 \cite{epw}. In rank one and two the degree bound makes $P_N$ constant; the highest coefficient of \eqref{rec} makes that constant $1$.

For $U_{3,4}$ the flats are the empty set, the four singletons, the six two-element subsets, and $E$. Their restrictions have characteristic polynomials, respectively,
\[
 1,\qquad t-1,\qquad (t-1)^2,\qquad
 \chi_M(t)=t^3-4t^2+6t-3.
\]
All nonempty-flat contractions have rank at most two, hence KL polynomial $1$. Write $P_M(t)=a+bt$. Equation \eqref{rec} becomes
\begin{align*}
 at^3+bt^2
  &=a+bt+4(t-1)+6(t-1)^2+(t^3-4t^2+6t-3)\\
  &=t^3+2t^2+(b-2)t+(a-1).
\end{align*}
Comparison of coefficients gives $a=1$ and $b=2$. Conversely, these values satisfy the whole identity. Thus
\[
 P_M(t)=1+2t,\qquad P_{M^*}(t)=1.
\]
The second equality follows from the rank-one calculation, or directly from $t=1+(t-1)$ for $U_{1,4}$.

\section{The divisibility contradiction}
The chromatic polynomial of $C_4$ is
\[
 \chi_G(t)=(t-1)^4+(t-1)=t^4-4t^3+6t^2-3t.
\]
In particular $\chi_G(-1)=14$. This can also be obtained from the subset inclusion-exclusion formula
\[
 \chi_G(t)=\sum_{S\subseteq E}(-1)^{|S|}t^{c(V,S)},
\]
which is the formula used in Lean. On the matroid-characteristic reading,
\[
 \chi_M(-1)=(-1)^3-4(-1)^2+6(-1)-3=-14,
\]
and its absolute value is again $14$.
\begin{theorem}
Conjecture 00000000579 is false.
\end{theorem}
\begin{proof}
For $G=C_4$, the claimed difference is $(1+2)-1=2$. There is no integer $z$ with $2=14z$. Thus the integer-multiple assertion already fails. Any stronger requirement on that integer multiplier also fails.
\end{proof}

\section{Formalization and semantic coverage}
The project imports only Lean's \texttt{Std}; it does not require Mathlib. The ground set has four edges, with subsets represented by \texttt{Fin 16} bit masks. The following bridges are included.
\begin{enumerate}
\item \texttt{adjacent} uses the explicit cycle edges above. \texttt{reachable} recursively computes walks of bounded length. Theorems \texttt{reach\_stabilizes}, \texttt{reach\_later}, and \texttt{connectivity\_exact} establish that four rounds equal connectivity by an arbitrary finite walk. Components are counted by their least vertex, including isolated vertices.
\item \texttt{graphic\_rank\_certificate} checks $4-c(V,S)=\min(3,|S|)$ for every subset. \texttt{graphicDualRank} is the actual dual formula; \texttt{graphic\_dual\_eq} identifies it with the verified rank table. The rank-zero, cardinality, monotonicity and submodularity axioms are checked for both tables.
\item For an interval $A\subseteq B$ with $A$ flat in $M|B$, \texttt{chi} computes the characteristic polynomial of $(M|B)/A$ as
\[
 \sum_{A\subseteq S\subseteq B}(-1)^{|S|-|A|}t^{r(B)-r(S)}.
\]
The ground set is $B\setminus A$: the formula introduces no extra contracted elements.
\item \texttt{coeff} represents every polynomial allowed by the strict degree bound at ranks zero through three. \texttt{IsKL} imposes rank-zero normalization and every coefficient, indexed by an arbitrary natural number, of \eqref{rec} for each loopless interval. At nonflat intervals unused data have no role in the recurrence.
\item \texttt{c0\_origin}, \texttt{c1\_origin}, and \texttt{dual\_c0} derive the relevant coefficients for \emph{any} family satisfying these axioms. They do not assume the desired polynomials. The theorems concern normalized coefficients, not the unused entries of the raw family.
\item \texttt{proposed\_isKL} and \texttt{dual\_proposed\_isKL} prove nonvacuity by constructing solutions on all relevant intervals. Finite checking handles coefficients zero through three. The general theorem \texttt{rhs\_tail} proves every higher coefficient vanishes, so no finite-degree cutoff is an assumption.
\item \texttt{chromatic\_coefficients}, \texttt{chromatic\_value}, and \texttt{characteristic\_coefficients} verify the inclusion-exclusion evaluations. \texttt{graphic\_counterexample} states the contradiction for the actual graph and dual-rank functions. \texttt{counterexample\_exists} includes existence of the recurrence families.
\end{enumerate}
All finite checks use kernel-reduced \texttt{decide}; symbolic steps use ordinary proof terms and \texttt{omega}. No \texttt{native\_decide}, proof holes, added axioms, external oracle, or opaque computational certificate is used. The printed dependencies of the main theorem are only the standard logical axioms \texttt{propext}, \texttt{Classical.choice}, and \texttt{Quot.sound}. The Python verifier is supplementary and is not trusted by Lean.

\section{Reproduction}
With Lean 4.31.0 available, run \texttt{lake build}. For direct verification run \texttt{lean Submission/Basic.lean}. Run \texttt{python3 verify.py} for a separate exhaustive arithmetic check. Rebuild this report with \texttt{pdflatex report.tex}. The included build logs record the actual checks performed for this submission.

\begin{thebibliography}{9}
\bibitem{statement} The Last Math Competition, conjecture 00000000579, retrieved 10 October 2026. \href{https://github.com/The-Last-Math-Competition/The-Last-Math-Competition/blob/main/conjectures/00000000579.md}{Original bilingual conjecture}.
\bibitem{epw} B. Elias, N. Proudfoot and M. Wakefield, \emph{The Kazhdan--Lusztig polynomial of a matroid}, Theorem 2.2, arXiv:1412.7408v3 (2016). \url{https://arxiv.org/abs/1412.7408}.
\end{thebibliography}
\end{document}
